\documentclass[11pt,a4paper]{article}

\usepackage[T1]{fontenc}
\usepackage[utf8]{inputenc}
\usepackage[brazil]{babel}
\usepackage{lmodern}
\usepackage{microtype}
\usepackage[margin=2cm]{geometry}
\usepackage{hyperref}
\usepackage{titlesec}
\usepackage{xcolor}

\hypersetup{
  colorlinks=true,
  linkcolor=black,
  urlcolor=black,
  citecolor=black,
  pdfborder={0 0 0}
}

\definecolor{sectioncolor}{HTML}{111111}
\definecolor{accentcolor}{HTML}{1A1A1A}

\pagestyle{empty}
\setlength{\parindent}{0pt}
\setlength{\parskip}{5pt}
\linespread{1.05}
\selectfont

\titleformat{\section}
  {\bfseries\large\color{sectioncolor}}
  {}{0pt}{}[\vspace{-2pt}{\color{sectioncolor}\titlerule[0.8pt]}]
\titlespacing*{\section}{0pt}{12pt}{8pt}

\newcommand{\cvrole}[4]{%
  \vspace{4pt}%
  {\bfseries\color{accentcolor}#1}\hfill{\itshape #2}\\[-2pt]%
  {\itshape #3}\hfill #4\\[6pt]%
}

\begin{document}

\begin{center}
  {\fontsize{20}{24}\selectfont\bfseries Kamille Hartmann Maccagnan}\\[8pt]
  {\large Governança de Dados \textbar{} Qualidade e Processos \textbar{} Estratégia de Dados}\\[10pt]
  Campo Comprido, Curitiba-PR \enspace\textbullet\enspace (41) 9 9928-4424 \enspace\textbullet\enspace
  \href{mailto:kamillehartmannm@gmail.com}{kamillehartmannm@gmail.com}\\[3pt]
  \href{https://linkedin.com/in/kamille-hartmann-maccagnan/}{linkedin.com/in/kamille-hartmann-maccagnan/}
\end{center}

\vspace{6pt}

\section{Resumo Profissional}

Profissional em formação na área de tecnologia, com experiência em governança de dados, qualidade de informações, documentação de processos e apoio a iniciativas de dados em ambiente corporativo e multinacional. Atua no acompanhamento de processos de governança, validação e padronização de dados, monitoramento de indicadores operacionais e interface entre áreas de negócio e T.I. Vivência com elaboração de artefatos de documentação, POPs, relatórios e dashboards para apoio à maturidade analítica e à tomada de decisão. Perfil analítico, organizado e colaborativo, com visão estratégica e interesse em elevar a confiabilidade dos dados como ativos organizacionais.

\section{Experiência Profissional}

\cvrole{Anjuss}{Analista de Suporte de TI}{07/2025 -- Atual}

Atuo na padronização de processos e na organização de informações operacionais, com elaboração e atualização de documentações e POPs que asseguram rastreabilidade, consistência e qualidade nas rotinas da rede de lojas e da equipe de T.I. Realizo validação de dados e testes manuais do sistema antes da liberação para as lojas, identificando inconsistências e reportando ocorrências para correção. Desenvolvo dashboards em Power BI para monitoramento de indicadores operacionais, apoiando a análise de desempenho e a identificação de oportunidades de melhoria nos fluxos. Colaboro com equipes internas para alinhar demandas entre usuários de negócio e áreas técnicas, garantindo comunicação clara e padronização das entregas.

\cvrole{HORSE do Brasil}{Estagiária de TI (IS/IT LATAM)}{07/2024 -- 07/2025}

Atuei em frentes de governança de dados e PMO em ambiente multinacional industrial, participando do desenvolvimento e acompanhamento de iniciativas de dados com foco em qualidade, documentação e confiabilidade das informações. Realizei verificação de qualidade de dados, organização de documentação de processos e apoio à execução de rotinas de governança para análises e relatórios confiáveis. Elaborei relatórios executivos e dashboards em Power BI para monitoramento de KPIs operacionais e de projetos, apoiando a tomada de decisão. Atuei como ponte entre áreas de negócio, operações, T.I. e stakeholders internacionais, participando de reuniões multilíngues e assegurando alinhamento entre as partes envolvidas. Utilizei o ServiceNow para gestão e acompanhamento de demandas relacionadas a dados e processos.

\cvrole{Restaurante Familiar}{Analista de Dados e Suporte Técnico}{01/2023 -- 07/2024}

Estruturei do zero a base de dados operacional e desenvolvi dashboards em Power BI para acompanhamento de indicadores de desempenho, com validação e padronização de informações para maior confiabilidade analítica. Realizei tratamento, integração e verificação de qualidade de dados de diferentes fontes operacionais, apoiando a consistência dos controles e relatórios gerenciais. Interagi com áreas operacionais e administrativas para garantir alinhamento entre necessidades de negócio e entregas analíticas, contribuindo para melhoria contínua dos processos de informação.

\section{Competências Técnicas}

\textbf{Governança de Dados:} Qualidade de dados, documentação de processos e informações, padronização de dados, validação e verificação de consistência, acompanhamento de processos de governança, monitoramento de indicadores, apoio a iniciativas de dados, melhoria contínua.

\textbf{Dados e Análise:} SQL, Power BI (DAX, Power Query), Excel avançado, ETL com Power Query, KPIs, relatórios gerenciais, análise de desvios, bancos de dados relacionais.

\textbf{Processos e Colaboração:} Interface entre negócio e T.I., documentação e POPs, gestão de demandas, comunicação com stakeholders, trabalho em equipe multidisciplinar, ServiceNow, Trello, Monday.

\textbf{Idiomas:} Inglês Avançado \enspace\textbullet\enspace Espanhol Básico \enspace\textbullet\enspace Coreano Básico

\section{Projetos e Conquistas}

1º lugar no Transformation Day Renault 2023 com solução orientada a dados e melhoria de processos. Participação no Hackathon Bosch 2023 com foco em resolução de problemas de negócio. Atuação em governança de dados e monitoramento de indicadores em ambiente multinacional na HORSE. Estruturação de base de dados e dashboards com validação de qualidade para apoio à gestão. Desenvolvimento de documentações e POPs para padronização e rastreabilidade de processos.

\section{Formação Acadêmica}

\textbf{PUCPR} \hfill Curitiba, PR\\
Graduação em Gestão da Tecnologia da Informação \hfill Previsão de conclusão: 06/2026\\[6pt]

\textbf{Escola Conquer} \hfill Curitiba, PR\\
Pós-graduação em Liderança e Gestão em Tecnologia \hfill Conclusão: 2024\\[6pt]

\textbf{Universidade Tuiuti do Paraná} \hfill Curitiba, PR\\
Graduação em Medicina Veterinária \hfill 06/2019\\[6pt]

\textbf{Qualittas} \hfill Curitiba, PR\\
Pós-graduação em Clínica Médica e Cirúrgica de Animais Exóticos e Pets Não Convencionais \hfill 12/2022

\end{document}
